% !TeX root = ../../tesi.tex

\section{Processi di sviluppo e gestione dei progetti}

\subsection{Organizzazione del lavoro e gestione dei progetti}
\label{sec:organizzazione-halue}

Il metodo di riferimento è \gloss{scrum}{Scrum}, un \emph{framework} agile per lo sviluppo
iterativo, adattato al lavoro di un’agenzia. Il lavoro procede per \gloss{sprint}{\emph{sprint}},
cicli brevi al termine dei quali il \emph{team} presenta i risultati al cliente in una \emph{sprint
review} e ne discute internamente in una retrospettiva. Mancano le figure canoniche di Scrum, il
\emph{Product Owner} e lo \emph{Scrum Master}, e le loro responsabilità passano a due ruoli:

\begin{itemize}
	\item il \gloss{team-lead}{\emph{team lead}}, che cura l’analisi funzionale, le stime per il
	cliente e il coordinamento tecnico;
	\item il \gloss{project-service-manager}{\emph{project/service manager}}, che pianifica le
	attività e la capacità del \emph{team}.
\end{itemize}

Accanto a loro lavorano gli sviluppatori e chi progetta l’interfaccia. In un’agenzia lo stesso
gruppo di persone lavora per più clienti nello stesso periodo, perciò nessuno è assegnato a un
progetto in modo stabile: ogni mese un \emph{demand planning} condiviso ripartisce il \emph{team}
sui clienti, e un rapporto sullo \gloss{sal}{stato di avanzamento dei lavori (SAL)} riporta con un
semaforo lo stato di ciascun progetto, le scadenze, i rischi e i punti aperti.

Il \emph{team} si coordina con gli strumenti della \tabref{tbl:coordinamento} e con uno
\gloss{stand-up}{\emph{stand-up}} quotidiano, una riunione mattutina di quindici-venti minuti in
cui ciascuno espone su cosa ha lavorato il giorno prima e su cosa lavorerà quel giorno.

\begin{tabellatesi}{P{0.3}P{0.7}}
	{Gli strumenti di coordinamento del \emph{team} di Halue.}
	{tbl:coordinamento}
	\toprule
	\rowcolor{TabellaIntestazione} \textbf{Strumento} & \textbf{Uso} \\
	\midrule
	Google Meet & riunioni a distanza, tra cui lo \emph{stand-up} \\
	Jira & registro delle richieste e \emph{board} di avanzamento dei progetti \\
	Slack & comunicazione con i clienti, in un canale per ciascun marchio \\
	\emph{Email} & resoconti delle attività e variazioni contrattuali \\
	Confluence e Google Drive & documentazione e conoscenza interna \\
	\bottomrule
\end{tabellatesi}

Gli orari e il luogo di lavoro sono flessibili: i dipendenti lavorano in sede indicativamente tra
il 30\% e il 40\% del tempo e da remoto per il resto, e ognuno gestisce la propria presenza in
autonomia. L’inserimento di un nuovo membro del \emph{team} parte dalla documentazione e dalle
procedure interne e prosegue con piccoli \gloss{ticket}{\emph{ticket}}, le unità di lavoro che Jira
registra con un identificativo, una descrizione e uno stato, su un progetto reale e con
l’affiancamento di un collega.

\subsection{Dal \emph{ticket} al rilascio}
\label{sec:ciclo-ticket}

La \figref{fig:ticket} mostra il percorso di una richiesta del cliente, dal \emph{ticket} su Jira
alla manutenzione.

\begin{figure}[htbp]
	\centering
	\setstretch{1}
	% Il ciclo di vita di un ticket in Halue, come diagramma di attività UML.
% Le partizioni sono le tre fasi; il flusso va da sinistra a destra su due righe.
% Il disegno è largo quanto il testo (150 mm): si include con \input, senza ridimensionarlo.

\begin{tikzpicture}[
	every node/.style={text=FiguraTesto},
	azione/.style={draw=FiguraTenue, fill=white, rounded corners=3pt, text width=18.7mm,
		minimum width=19.8mm, minimum height=15mm, align=center, inner sep=1.5pt,
		font=\fontsize{7.5}{9}\selectfont,
		execute at begin node={\hyphenpenalty=10000\exhyphenpenalty=10000}},
	decisione/.style={draw=FiguraTenue, fill=white, diamond, minimum size=6mm, inner sep=0pt},
	iniziale/.style={circle, fill=FiguraTesto, minimum size=3mm, inner sep=0pt},
	finale/.style={circle, draw=FiguraTesto, fill=white, minimum size=3.6mm, inner sep=0pt,
		path picture={\fill[FiguraTesto] (path picture bounding box.center) circle (1.05mm);}},
	flusso/.style={freccia},
	guardia/.style={font=\fontsize{7.5}{9}\selectfont, text=FiguraTenue, inner sep=1pt},
	node distance=4mm,
]
	% Prima riga: analisi.
	\node[iniziale] (inizio) at (4mm, -20mm) {};
	\node[azione, right=of inizio] (richiesta)
		{\voce{Richiesta}{del cliente per \emph{call}, \emph{email} o \emph{workshop}}};
	\node[azione, right=of richiesta] (ticket)
		{\voce{\emph{Ticket} su Jira}{\emph{Epic}, \emph{Story}, \emph{Change Request} o \emph{Bug}}};
	\node[azione, right=of ticket] (funzionale)
		{\voce{\emph{Refinement} funzionale}{completezza e fattibilità}};
	\node[decisione, right=of funzionale] (bug) {};
	\node[azione, right=7mm of bug] (tecnico)
		{\voce{\emph{Refinement} tecnico}{svolto dagli sviluppatori}};
	\node[decisione, right=of tecnico] (unione) {};
	\node[azione, right=of unione] (stima)
		{\voce{Stima}{\emph{story point} per il \emph{team}, ore per il cliente}};

	% Seconda riga: sviluppo, verifica e rilascio.
	\node[azione] (branch) at (12.4mm, -52mm)
		{\voce{Sviluppo}{su un \emph{branch} con la chiave del \emph{ticket}}};
	\node[azione, right=3.5mm of branch] (pr)
		{\voce{\emph{Pull request}}{\emph{test} automatici e manuali}};
	\node[azione, right=3.5mm of pr] (designer)
		{\voce{Revisione}{del \emph{designer} prima del rilascio}};
	\node[azione, right=5mm of designer] (collaudo)
		{\voce{Collaudo}{del cliente in \emph{staging}}};
	\node[azione, right=3.5mm of collaudo] (produzione)
		{\voce{Rilascio in produzione}{\emph{deploy} avviato a mano}};
	\node[azione, right=3.5mm of produzione] (manutenzione)
		{\voce{Manutenzione}{le segnalazioni diventano nuovi \emph{ticket}}};
	\node[finale, right=3.5mm of manutenzione] (fine) {};

	% Riferimenti per partizioni e percorsi a gomito.
	\coordinate (scavalco) at (0, -9.5mm);
	\coordinate (ritorno) at (0, -33mm);
	\coordinate (riga2) at (0, -36mm);
	\coordinate (fondo2) at (0, -62mm);
	\coordinate (confine) at ($(designer.east)!0.5!(collaudo.west)$);
	\coordinate (confine-alto) at (confine |- riga2);
	\coordinate (confine-basso) at (confine |- fondo2);

	% Partizioni, sullo sfondo, con l'intestazione nello stile delle tabelle.
	\begin{scope}[on background layer]
		\path[partizione] (0, 0) rectangle (150mm, -30mm);
		\path[intestazione] (0, 0) rectangle (150mm, -6mm);
		\node[titolo] at (75mm, -3mm) {Analisi};
		\path[partizione] (riga2) rectangle (confine-basso);
		\path[intestazione] (riga2) rectangle ($(confine-alto) + (0, -6mm)$);
		\node[titolo] at ($(riga2)!0.5!(confine-alto) + (0, -3mm)$) {Sviluppo};
		\path[partizione] (confine-alto) rectangle (150mm, -62mm);
		\path[intestazione] (confine-alto) rectangle (150mm, -42mm);
		\node[titolo] at ($(confine-alto)!0.5!(150mm, -36mm) + (0, -3mm)$) {Verifica e rilascio};
	\end{scope}

	% Flusso di controllo.
	\draw[flusso] (inizio) -- (richiesta);
	\draw[flusso] (richiesta) -- (ticket);
	\draw[flusso] (ticket) -- (funzionale);
	\draw[flusso] (funzionale) -- (bug);
	\draw[flusso] (bug) -- node[guardia, above] {[no]} (tecnico);
	\draw[flusso] (bug.north) -- (bug.north |- scavalco)
		-- node[guardia, above] {[sì]} (unione.north |- scavalco) -- (unione.north);
	\draw[flusso] (tecnico) -- (unione);
	\draw[flusso] (unione) -- (stima);
	\node[guardia, below=1mm of bug] {È un \emph{bug}?};
	\draw[flusso] (stima.south) -- (stima.south |- ritorno) -- (branch.north |- ritorno)
		-- (branch.north);
	\draw[flusso] (branch) -- (pr);
	\draw[flusso] (pr) -- (designer);
	\draw[flusso] (designer) -- (collaudo);
	\draw[flusso] (collaudo) -- (produzione);
	\draw[flusso] (produzione) -- (manutenzione);
	\draw[flusso] (manutenzione) -- (fine);
\end{tikzpicture}

	\caption{Il ciclo di vita di un \emph{ticket} in Halue, come diagramma di attività UML.}
	\label{fig:ticket}
\end{figure}

Gli \emph{Epic} raccolgono insiemi di funzionalità, le \emph{Story} descrivono singole
funzionalità, le \emph{Change Request} le modifiche chieste a progetto avviato e i \emph{Bug} i
difetti.

La \tabref{tbl:pratiche} raccoglie, nell’ordine della \figref{fig:ticket}, le pratiche che
accompagnano il ciclo. Il metodo è rigido dove la struttura si ripaga sempre, cioè nella
tracciabilità dalla richiesta al codice e nelle stime, e flessibile dove i clienti sono diversi tra
loro, come la copertura dei \emph{test} o i livelli di servizio, che si decidono progetto per
progetto. La stessa impostazione ha guidato il mio \emph{stage}, per esempio nella scelta di
concentrare i \emph{test} automatici solo dove servivano (\secref{sec:verifica}).

\begin{tabellatesi}{P{0.24}P{0.76}}
	{Le pratiche di Halue lungo il ciclo di un \emph{ticket}.}
	{tbl:pratiche}
	\toprule
	\rowcolor{TabellaIntestazione} \textbf{Ambito} & \textbf{Pratica} \\
	\midrule
	Stima per il \emph{team} & \gloss{story-point}{\emph{story point}}, una misura della complessità
	relativa sui numeri di Fibonacci da 1 a 13; oltre 13 il \emph{ticket} si divide \\
	Stima per il cliente & ore, con un metodo interno che pesa complessità, chiarezza della
	richiesta, \emph{team} coinvolto e confidenza nella stima \\
	Decisioni di architettura & scritte nel \emph{ticket}, oppure in un \gloss{adr}{ADR}, un
	documento breve con contesto, scelta e motivazione, se valgono per tutto il progetto \\
	Progetto grafico & da Figma, collegato al \emph{ticket}, in versione per dispositivi mobili e per
	\emph{desktop} \\
	Codice & una \gloss{pull-request}{\emph{pull request}}, cioè una proposta di integrazione delle
	modifiche, per ogni \emph{ticket}, con la sua chiave nei messaggi di \emph{commit}; stile
	controllato da ESLint e Prettier \\
	Versioni & \gloss{semver}{versionamento semantico (SemVer)}, che distingue modifiche
	incompatibili, nuove funzionalità e correzioni, con \emph{tag} e \emph{release} su GitHub \\
	\emph{Test} & nessuna copertura minima prefissata \\
	Collaudo & in \gloss{staging}{\emph{staging}} o \emph{preview}, un ambiente separato che
	riproduce le condizioni di produzione \\
	\emph{Deploy} & automatizzato, secondo una procedura scritta \\
	Manutenzione & segnalazioni da Jira, Slack o \emph{email}; priorità concordate caso per caso,
	senza livelli di servizio formali \\
	Retrospettive & anche alla chiusura di un progetto o di una \emph{release} importante \\
	\bottomrule
\end{tabellatesi}
